\section*{Chapitre \ref{ch:b3:computational-chemistry} --- Chimie numérique~: Hartree--Fock et DFT}

\begin{solution}{exo:b3:computational-chemistry:1}
$\dd E/\dd\alpha = \frac32 - \sqrt{2/\pi\alpha} = 0$ donne $\alpha = 8/9\pi = 0.283\,
a_0^{-2}$ et $E = \frac32\alpha - 2\sqrt{2\alpha/\pi} = -4/3\pi = \qty{-0.4244}{\hartree}$,
15\,\% au-dessus de la valeur exacte $-0.5$~: une gaussienne seule n'a ni
le point de rebroussement ni la queue de la fonction $1s$.
\end{solution}

\begin{solution}{exo:b3:computational-chemistry:2}
STO-3G~: 5 fonctions par C ($1s$, $2s$, trois $2p$) et 1 par H~: $6 \times 5 + 6 =
36$. 6-31G(d)~: par C, 1 (cœur) + $2 \times 4$ (valence séparée) + 6
($d$) $= 15$~;
par H, 2~: $6 \times 15 + 6 \times 2 = 102$.
\end{solution}

\begin{solution}{exo:b3:computational-chemistry:3}
Le hamiltonien électronique ne contient pas les masses des noyaux~: dans
l'\omterm{def:b3:computational-chemistry:born-oppenheimer}{approximation de Born--Oppenheimer}, la SEP, donc $r_e$ et
$k = U''(r_e)$, est la
même. $\tilde\omega = \sqrt{k/\mu}/2\pi c$ dépend de la \omterm{def:b3:quantum-model-systems:oscillator}{masse réduite}, qui
double~: rapport $\sqrt2 = 1.414$ (1.413 mesuré).
\end{solution}

\begin{solution}{exo:b3:computational-chemistry:4}
Huit atomes donnent $3N - 6 = 18$ vibrations~: 17 réelles et une
imaginaire. Une \omterm{def:b3:quantum-model-systems:operator}{valeur propre} négative de la hessienne~: un point selle
d'ordre un, une structure de transition (ici, par exemple, d'une rotation
ou d'une élimination).
\end{solution}

\begin{solution}{exo:b3:computational-chemistry:5}
$\zeta = 27/16$, $E = -(27/16)^2 = \qty{-2.8477}{\hartree} = \qty{-77.49}{eV}$.
Erreur $-2.8477 + 2.9034 = \qty{0.0557}{\hartree} = \qty{1.52}{eV}$, soit
1.9\,\% de l'énergie totale --- mais davantage que bien des énergies de
réaction.
\end{solution}

\begin{solution}{exo:b3:computational-chemistry:6}
$(-1 - E)(-0.5 - E) - (-0.2 - 0.3E)^2 = 0$, c'est-à-dire $0.91E^2 + 1.38E + 0.46 = 0$~:
$E = \qty{-1.0217}{\hartree}$ et \qty{-0.4947}{\hartree}. Oui~: le
mélange avec la seconde fonction abaisse l'énergie au-dessous de
$H_{11}$, comme le permet le \omterm{thm:b3:computational-chemistry:variational}{principe variationnel}.
\end{solution}

\begin{solution}{exo:b3:computational-chemistry:7}
$(102/36)^4 = 64$~: environ 640 minutes, soit quelque 11 heures.
\end{solution}

\begin{solution}{exo:b3:computational-chemistry:8}
$0.5782 \times 27.211 = \qty{15.73}{eV}$, 0.30 eV au-dessus de la valeur
mesurée. Orbitales figées~: l'ion réel se relaxe et se trouve plus bas,
donc la valeur de Koopmans est trop haute. Corrélation manquante~: la
molécule neutre (deux électrons) a plus d'\omterm{def:b3:computational-chemistry:correlation}{énergie de corrélation} que
l'ion (un électron), ce qui rend la vraie valeur plus grande. Ici, la
première erreur l'emporte.
\end{solution}

\begin{solution}{exo:b3:computational-chemistry:9}
$-0.5960 - 2(-0.4666) = \qty{0.3372}{\hartree} = \qty{9.18}{eV}$ au-dessus
de deux atomes. La fonction RHF conserve 50\,\% de \ce{H+ H-}, et séparer
un proton et un hydrure coûte la différence entre l'énergie d'ionisation
de H et l'affinité électronique de H, que l'énergie inclut dans sa
moyenne.
\end{solution}

\begin{solution}{exo:b3:computational-chemistry:10}
$E(c_1^2 + 2c_1c_2S + c_2^2) = c_1^2H_{11} + 2c_1c_2H_{12} + c_2^2H_{22}$.
En dérivant par rapport à $c_1$ avec $\partial E/\partial c_1 = 0$~:
$E(2c_1 + 2c_2S) = 2c_1H_{11} + 2c_2H_{12}$, c'est-à-dire $(H_{11} - E)c_1 + (H_{12} -
ES)c_2 = 0$~; de même $(H_{12} - ES)c_1 + (H_{22} - E)c_2 = 0$.
\end{solution}

\begin{solution}{exo:b3:computational-chemistry:11}
$k \approx [E(1.30) - 2E(1.35) + E(1.40)]/(0.05)^2 = 0.001417/0.0025 =
\qty{0.567}{\hartree\per\bohr\squared} = \qty{882}{N/m}$. Avec $\mu = m_H/2 =
\qty{8.37e-28}{kg}$, $\omega = \sqrt{k/\mu} = \qty{1.03e15}{s^{-1}}$ et
$\tilde\omega = \omega/2\pi c = \qty{5452}{cm^{-1}}$, 24\,\% au-dessus de
la valeur mesurée
$\tilde\omega_e$~: une courbe RHF en \omterm{def:b3:computational-chemistry:basis}{base minimale} est trop raide (pas de
corrélation, base trop petite), et c'est pourquoi on applique un facteur
d'échelle inférieur à 1 aux fréquences calculées.
\end{solution}

\begin{solution}{exo:b3:computational-chemistry:12}
Des écarts de quelques \unit{kJ/mol} mettant en jeu une liaison hydrogène
exigent la corrélation et la dispersion~: une fonctionnelle hybride avec
correction de dispersion, ou mieux une méthode corrélée de fonction
d'onde pour les énergies finales, avec une base polarisée triple zêta
comportant des \omterm{def:b3:computational-chemistry:split-valence}{fonctions diffuses}. Optimiser les deux conformères au même
niveau, confirmer les minima par les fréquences (toutes réelles), ajouter
les \omterm{def:b3:quantum-model-systems:oscillator}{énergies de point zéro}, tester la base par un calcul plus grand, et
comparer à un système apparenté d'énergie conformationnelle connue.
\end{solution}

\begin{solution}{pb:b3:computational-chemistry:1}
\textbf{1.} Chaque \omterm{def:b3:computational-chemistry:basis}{orbitale de type Slater} d'une \omterm{def:b3:computational-chemistry:basis}{base minimale} est
remplacée par une contraction fixe de trois gaussiennes ajustées sur
elle.
\textbf{2.} Le noyau d'hélium attire davantage ses électrons, donc sa
fonction $1s$ est plus compacte (exposants plus grands). $6.3624/3.4253 = 1.857 = (1.69/1.24)^2$.
\textbf{3.} Deux fonctions de base, deux électrons, une orbitale occupée
(et une virtuelle).
\textbf{4.} $V_{\mathrm{nn}} = Z_{\mathrm{He}}Z_{\mathrm H}/R = 2/1.4632 =
\qty{1.3669}{\hartree}$.
\textbf{5.} Les deux fonctions se recouvrent fortement (54\,\%)~: les
atomes sont assez proches pour se lier.
\textbf{6.} $h_{11}$ est l'énergie d'un électron dans la fonction de He
en présence des deux noyaux, mais sans autre électron~; le noyau He
(charge 2) le retient bien plus fortement.
\textbf{7.} $(h_{11} - \varepsilon)(h_{22} - \varepsilon) - (h_{12} - \varepsilon S)^2 = 0$~:
$0.71185\varepsilon^2 + 2.79292\varepsilon + 2.44969 = 0$.
\textbf{8.} $\varepsilon = \qty{-2.600}{\hartree}$ et \qty{-1.324}{\hartree}.
\textbf{9.} $\mathbf h$ ignore la répulsion entre les deux électrons~;
l'orbitale de cœur est trop contractée et trop basse.
\textbf{10.} $F_{\mu\nu} = h_{\mu\nu} + \sum_{\lambda\sigma}P_{\lambda\sigma}[(\mu\nu|\sigma\lambda) -
\frac12(\mu\lambda|\sigma\nu)]$, avec $P_{\lambda\sigma} = 2C_{\lambda1}C_{\sigma1}$.
\textbf{11.} La répulsion coulombienne de chaque électron par le nuage de
charge de l'autre (et, en général, l'échange).
\textbf{12.} $\mathbf F$ dépend de $\mathbf P$, qui dépend des orbitales
obtenues à partir de $\mathbf F$~: les équations sont non linéaires et se
résolvent par approximations successives jusqu'à l'autocohérence.
\textbf{13.} $c_1^2 + c_2^2 + 2c_1c_2S = 0.7684 + 0.0410 + 0.1906 = 1.0000$.
\textbf{14.} He~: $1.5369 + 0.1906 = 1.727$~; H~: $0.0820 + 0.1906 = 0.273$
électron.
\textbf{15.} Charges~: He $+0.27$, H $+0.73$. La charge positive se trouve
surtout sur l'hydrogène~: \ce{HeH+} est un proton lié à un atome
d'hélium, \ce{He + H+}, comme on s'y attend puisque l'énergie
d'ionisation de He (24.6 eV) dépasse de loin celle de H (13.6 eV).
\textbf{16.} $-\varepsilon_1 = \qty{1.633}{\hartree} = \qty{44.4}{eV}$.
\textbf{17.} $E_{\mathrm{el}} = -2.8418 - 1.3669 = \qty{-4.2087}{\hartree}$.
\textbf{18.} Trois cycles atteignent \qty{-2.8418}{\hartree}. Chaque cycle
donne un déterminant, dont l'énergie majore l'énergie Hartree--Fock
convergée (\omterm{thm:b3:computational-chemistry:variational}{principe variationnel}), et les itérations l'améliorent.
\textbf{19.} La base avec $\zeta = 2.0925$ donne l'énergie la plus basse,
c'est donc la meilleure pour cette molécule (\omterm{thm:b3:computational-chemistry:variational}{principe variationnel})~: la
fonction de l'hélium est plus compacte dans \ce{HeH+} que dans l'atome
libre, dont la base standard reprend le $\zeta = 1.69$.
\textbf{20.} L'orbitale antiliante vacante $\sigma^*$~; son énergie
approcherait l'opposé de l'affinité électronique de \ce{HeH+} (mal, dans
une base aussi petite).
\textbf{21.} Zéro (aucun électron)~; $\qty{-2.8078}{\hartree}$.
\textbf{22.} $-2.8078 - (-2.8418) = \qty{0.0340}{\hartree} = \qty{89}{kJ/mol}$.
\textbf{23.} Environ la moitié de la valeur mesurée~: une \omterm{def:b3:computational-chemistry:basis}{base minimale} ne
peut pas polariser la densité de l'hélium vers le proton (il n'y a pas
de fonctions $p$), si bien que la liaison est trop faible. (Les
corrections de point zéro et à \qty{298}{K} sont petites à côté.)
\textbf{24.} $-(24.5874 + 54.4178)/27.211 = \qty{-2.9034}{\hartree}$~;
l'atome en STO-3G est trop haut de \qty{0.0956}{\hartree} (\qty{2.6}{eV}).
\textbf{25.} Les géométries dépendent de la variation de l'énergie avec
$R$~; l'essentiel de l'erreur, concentré dans le cœur, est presque le
même à toute géométrie et se compense.
\textbf{26.} \textbf{$E(\ce{HeH+}) = \qty{-2.8418}{\hartree}$ en RHF/STO-3G
(base standard), à \qty{1.4632}{\bohr}} (et \qty{-2.8607}{\hartree} avec
le $\zeta_{\mathrm{He}} = 2.0925$ classique).
\end{solution}
